\documentclass[11pt]{article}
\usepackage[T1]{fontenc}
\usepackage{lmodern,amsmath,amssymb,amsthm,booktabs,microtype}
\usepackage[a4paper,margin=26mm]{geometry}
\usepackage[hidelinks]{hyperref}
\usepackage{fancyhdr}
\newtheorem{theorem}{Theorem}
\newtheorem{proposition}[theorem]{Proposition}
\newtheorem{lemma}[theorem]{Lemma}
\newtheorem{corollary}[theorem]{Corollary}
\theoremstyle{remark}\newtheorem*{remark}{Remark}
\newcommand{\PP}{\mathbb P}
\newcommand{\floor}[1]{\lfloor #1\rfloor}
\newcommand{\dimH}{\dim_{\mathrm H}}
\setlength{\headheight}{14pt}\setlength{\emergencystretch}{2em}\raggedbottom
\pagestyle{fancy}\fancyhf{}\lhead{\small Prime outputs on two shifted Mills sequences}
\rhead{\small Research note}\cfoot{\thepage}
\hypersetup{pdftitle={Prime outputs on the two shifted Mills sequences 3 to the k plus or minus k}}
\title{Prime outputs on the two shifted Mills sequences\\$3^k-k$ and $3^k+k$}
\author{Research note prepared for Keio}
\date{3 November 2026}
\begin{document}
\maketitle
\begin{abstract}
The least prime-representing constants for $3^k-k$ and $3^k+k$
are separately transcendental in the corpus. We study the primes
and the surrounding integers on these sequences. A single suitably chosen base realizes any prescribed sequence of signs. At each
sufficiently late level the chosen floor is prime and every integer
within distance $ck$ of the opposite floor is composite. The set of bases has full local Hausdorff dimension. In every sufficiently
large dyadic interval, the first four levels can have both floors prime and
the prescribed choices can start at level five. For the two original
minima we also exhibit fixed real dilations whose new prime outputs
contain earlier primes from both sequences as a decimal suffix.
Other dilations prescribe small and large gaps to the next prime.
Finally, a three-term ratio of the original prime outputs recovers
a transcendental number and separates all competing monomials in
three consecutive terms, with an explicit obstruction at four terms.
\end{abstract}

\section{Which parameters are fixed?}
For $\sigma\in\{-1,1\}$, define
\[
 C_k^\sigma=3^k+\sigma k,\qquad
 A_\sigma=\min\{a>1:\lfloor a^{C_k^\sigma}\rfloor\in\PP
                  \text{ for all }k\ge1\},\qquad
 p_k^\sigma=\lfloor A_\sigma^{C_k^\sigma}\rfloor .
\]
Here $\PP$ is the set of ordinary positive primes. Manuscript~008 proves
existence and transcendence of the two minima as special cases of its
coefficient-family theorem \cite{corpus8}. Its proof verifies the
recurrence conditions in Saito's Theorem~2.3 \cite{saito}. The same
manuscript explicitly leaves existence of simultaneous prime floors on a common base unresolved.
The corresponding increasing exponents begin with $2,4,7,11,24,30,77,85$. Manuscript~031 also considers prime-free windows at powers of prescribed
Pisot or Salem numbers \cite{corpus31}.

We use three choices of parameters. In Section~\ref{sec:switch}, we
choose one base $a$ after specifying which sign is to produce a prime
at each level. In Sections~\ref{sec:suffix} and~\ref{sec:gaps},
the two original bases stay fixed. We use a new scalar $\lambda$
to act on their prime outputs. In the final section even this
scalar is absent. Only the last conclusion uses transcendence itself.
The others use the lengths of the two exponent sequences or the
rate at which their prime outputs grow.

For quick reference, put $P_k=\max(p_k^-,p_k^+)$ and
$L=\max(\log A_-,\log A_+)>0$, using $+$ and $-$ as superscripts for the
two signs. Rounding has exponentially small relative error, so
\begin{equation}\label{eq:growth}
 \log P_k=L3^k+O(k),\qquad
 \frac{\log P_{k+1}}{\log P_k}\longrightarrow3 .
\end{equation}
These statements require no ordering of the two minima.

\section{Selecting the prime floor while excluding nearby primes}
\label{sec:switch}

\begin{theorem}\label{th:switch}Put $d_k=3^k-k$ and $e_k=3^k+k$. Fix a bounded open interval
$J=(a,b)$ with $1<a<b$. There exist $c_J>0$ and an integer $K_J$
with the following property. For every sequence of signs
$(s_k)_{k\ge K_J}$, the set of $A\in J$ satisfying the following
conditions at every $k\ge K_J$ has Hausdorff dimension one.
\begin{itemize}
\item If $s_k=-$, then $\lfloor A^{d_k}\rfloor$ is prime and
      $\lfloor A^{e_k}\rfloor+h$ is composite for every integer
      $h$ with $|h|\le\lfloor c_J k\rfloor$.
\item If $s_k=+$, then $\lfloor A^{e_k}\rfloor$ is prime and
      $\lfloor A^{d_k}\rfloor+h$ is composite throughout the same range.
\end{itemize}
The quantifiers allow the signs to be specified before choosing the
base.
\end{theorem}

\begin{proof}
We use two classical prime estimates. Uniformly for $x^{3/5}\le h\le x$,
the interval $[x,x+h]$ contains $\gg h/\log x$ primes for all sufficiently
large $x$, by the prime number theorem in short intervals with exponent
$7/12+\varepsilon<3/5$ \cite{tz}. We also use the Brun--Titchmarsh inequality
in the form $\#\{p\text{ prime}:u<p\le u+h\}\ll h/\log h$ for $h\ge3$,
uniformly in $u\ge0$ \cite{mv}. Fixed changes of endpoints are harmless below.

At level $k-1$ every parent is a closed interval $P\subset J$ of the form
\[
 P=\{A:v+1/4\le A^{e_{k-1}}\le v+3/4\}
\]
for an integer $v$. In each parent retain children of this same form at
level $k$. Write $I_-=P^{d_k}$ and $I_+=P^{e_k}$ for the images of
$P$ under these increasing powers, $X,Y$ for their left endpoints,
and $H_-,H_+$ for their lengths. With $A_0=\inf P$ set
\[
 r=e_k/d_k,\qquad \lambda=r A_0^{2k}.
\]
Uniformly in the parent, differentiation of the powers gives
\[
 H_-\asymp X^{1-e_{k-1}/d_k},\qquad
 H_+\asymp Y^{1-e_{k-1}/e_k},\qquad H_+\asymp\lambda H_-.
\]
Indeed the map from $I_-$ to $I_+$ is $x\mapsto x^r$ with derivative
$rA^{2k}$, and its ratio to $\lambda$ tends to one uniformly on $P$,
since $A^{e_{k-1}}$ ranges over an interval of length $1/2$ at
$v\ge a^{e_{k-1}}-1$. The powers from the parent coordinate to
either $I_\pm$ have exponents tending to three. Consequently, for
all sufficiently large $k$, middle subintervals of either $I_\pm$
of any fixed positive fraction of its length are in the range of
the short-interval lower bound above. We have
$\log H_-\asymp\log X\asymp\log Y$ and
$\log\lambda\ge 2k\log a$.

Write $B=\lfloor c_J k\rfloor$. First suppose $s_k=+$. There are
$\gg H_+/\log Y$ primes $q$ for which the full unit cell $[q,q+1]$
lies strictly inside $I_+$. Remove a prime $q$ if its cell meets
the image under $x\mapsto x^r$ of any lower unit cell $[j,j+1]$
with $|j-p|\le B$ for some prime $p$. It suffices to consider primes
$p$ in $I_-$ enlarged by $B+2$ at both ends. Their number is
$\ll H_-/\log X$ by Brun--Titchmarsh, because $B=o(H_-)$.
For each such $p$, the relevant lower cells form an interval of
length at most $2B+1$, clipped to $I_-$. Its upper image, thickened
by one at both ends, has length at most $C(B+1)\lambda$.
Brun--Titchmarsh bounds its prime count by
$O((B+1)\lambda/\log\lambda)$, also when the image is shorter
(enlarge it to the stated upper length). The total number removed is
therefore
\[
 \ll \frac{H_+}{\log Y}\frac{B+1}{\log\lambda}.
\]
All implied constants here are uniform in $k$ and in the parent.
Choose $c_J>0$ small enough, and then $K_J$ large enough, that
at most half of the lower bound on the available $q$ is removed.
The central half of each surviving $q$-cell is a valid child. In
particular, a child may cross a boundary between two lower integer
cells, since both lower labels then satisfy the required exclusions.

Suppose instead $s_k=-$. There are $\gg H_-/\log X$ lower primes
$p$ with $[p,p+1]$ inside $I_-$. The upper image of
$[p+1/4,p+3/4]$ has length comparable to $\lambda$, so contains
$\gg\lambda$ full integer cells $[q,q+1]$. The number of primes in
this upper interval enlarged by $B+2$ at both ends is
$O(\lambda/\log\lambda)$, since $B=o(\lambda)$. Each prime rules out
at most $2B+1$ integer labels $q$. Decreasing $c_J$ and increasing
$K_J$ once for both cases leaves $\gg\lambda$ choices per lower
prime. Cells obtained from distinct $p$ are disjoint. All integers
subject to exclusion are greater than one once $k\ge K_J$.

For either sign the number of retained children is thus
$\gg H_+/\log Y$. Choose a parent of the stated form at level
$K_J-1$ strictly inside $J$, increasing $K_J$ if necessary. Repeating
the construction gives a compact nested limit set. To compute its
dimension, distribute the mass of each parent equally among its
retained children. For each child $Q$,
$|Q|/|P|\asymp1/H_+$ by the same derivative estimate, and hence
\[
 \frac{m(Q)}{m(P)}\ll(\log Y)\frac{|Q|}{|P|}.
\]
As $\log Y\le e_k\log b$, iteration gives
$m(Q)\le |Q|\exp(O_J(k^2))$ for every level-$k$ cell $Q$ (the
fixed starting interval is absorbed in the constant).
Given an interval of small length $\delta$, take
$k=O_J(\log\log(1/\delta))$ sufficiently large that every
level-$k$ cell has length at most $\delta$. This is possible
since $|Q|\le 1/(2e_k a^{e_k-1})$. The disjoint level-$k$
cells meeting the interval have combined length at most $3\delta$.
Their combined mass is at most
$3\delta\exp(O_J((\log\log(1/\delta))^2))=O_{s,J}(\delta^s)$
for every $s<1$. The mass distribution principle proves the claim.
\end{proof}

\paragraph{Lebesgue measure.}
The full Hausdorff dimension in Theorem~\ref{th:switch} coexists with
Lebesgue measure zero. In fact, for almost every $A\in J$, both floors
are composite at every sufficiently large level. For a positive integer
$E$, the substitution $x=A^E$ gives
\[
 \big|\{A\in(a,b):\lfloor A^E\rfloor\in\PP\}\big|
 \le \frac{b}{E}\!
       \sum_{\substack{a^E-1\le p\le b^E\\p\in\PP}}\frac1p
 \ll_{a,b} E^{-1} .
\]
The last bound follows by partial summation from $\pi(x)\ll x/\log x$.
Both sums $\sum_k 1/C_k^\sigma$ converge. The first Borel--Cantelli
lemma therefore gives the assertion without any independence hypothesis.

\subsection{Four simultaneous levels before the prescribed tail}
\label{sec:fourdouble}

\begin{corollary}\label{cor:four}For every sufficiently large $R$ and every prescribed sequence
$(s_k)_{k\ge5}$, a set of bases of Hausdorff dimension one in $(R,2R)$
has both floors prime at each $1\le k\le4$ and has the prime and
composite choices of Section~\ref{sec:switch} at every $k\ge5$.
The composite buffer in this assertion may be taken to include every
integer shift $|h|\le k$.
\end{corollary}

\begin{proof}
For $t$ a positive integer, call the inverse image of
$[p+1/4,p+3/4]$ under $A\mapsto A^t$ a trimmed prime cell if $p$ is
prime. We first prove that the union of exponent-$85$ trimmed prime
cells contained in successive trimmed prime cells at the seven earlier
exponents in the list
\[
        2,\ 4,\ 7,\ 11,\ 24,\ 30,\ 77,\ 85
\]
has measure $\gg R/(\log R)^8$ inside $(R,2R)$. These are precisely
the two exponents at each of the first four levels.

We use the mean-square error term in Watt's Theorem~1 \cite{watt}. For fixed $D$ and $\varepsilon>0$,
$X^{1/14+\varepsilon}\le Y<X/2$, and $x\in[X/2,X]$,
put $y=xY/X$. Except on a set of measure $O_{D,\varepsilon}
(X/(\log X)^D)$, the interval $(x-y,x]$ contains $\gg Y/\log X$
primes. Indeed, Watt's error in normalized $L^2$ is
$O_E(Y/(\log X)^E)$, so Chebyshev's inequality gives this statement
on taking $2E-2\ge D$.

The prime number theorem applied at scale $R^2$ yields a union of
exponent-$2$ trimmed prime cells of total length $\gg R/\log R$.
Suppose the parent cells $P$ at one exponent $m$ have total measure
$\gg R/(\log R)^j$ after $j$ steps, and let $n$ be the next exponent.
Their images $P^n=\{A^n:A\in P\}$ are disjoint and have lengths
$\asymp_{m,n}R^{n-m}$. Fix $Y=\eta_{m,n} R^{n-m}$ with
$\eta_{m,n}>0$ small enough that $2Y$ is shorter than one quarter of
every such image. Cover $[R^n,(2R)^n]$ by $O_n(1)$ dyadic intervals
$[X/2,X]$ and apply the preceding estimate with $D=12$ in each.
The bound is applicable since, for each of the seven transitions,
\[
              1-m/n\ \ge\ 8/85\ >\ 1/14 .
\]
Across all these intervals the union of exceptional sets has measure
$O_{m,n}(R^n/(\log R)^{12})$. Call a parent bad if the middle third
of $P^n$ lies wholly in this union. Because $(A^n)'\asymp_n R^{n-1}$
on $(R,2R)$, all bad parents combined have measure
$O_{m,n}(R/(\log R)^{12})$. In each other parent choose a
nonexceptional $x$ in that middle third and use its dyadic interval.
Then $(x-y,x]$ lies inside $P^n$ with room for full unit cells around
all the counted primes. The corresponding trimmed exponent-$n$
children have total measure $\gg_{m,n}|P|/\log R$.
Discarding the bad parents still leaves $\gg R/(\log R)^j$ total
parent length. Induction through $j=8$ proves the asserted measure bound.

Choose any one of the resulting nonempty exponent-$85$ cells as the
initial parent for the proof of Section~\ref{sec:switch}. Its
short-interval and counting estimates apply uniformly from $k=5$ on
when $R$ is sufficiently large. Specifically, for $k\ge5$,
$e_{k-1}/d_k\le 3/8$ (equivalently $11k-8\le 3^{k-1}$).
The ratios $d_k/e_{k-1}$ and $e_k/e_{k-1}$ stay in a fixed compact
subinterval of $(1,4)$, so the derivative comparisons there have
absolute constants. Hence $H_-\gg X^{5/8}$ and $H_+\gg Y_+^{5/8}$,
where $Y_+$ denotes the left endpoint of the upper image.
With $B=k$, the proportion lost to each exclusion estimate is bounded
by an absolute constant times
$(k+1)/(2k\log R)$, uniformly less than the required fixed threshold
as $R$ grows. The mass estimate from that section now yields dimension
one in this starting cell.
\end{proof}

\subsection{A rational base with simultaneous prime floors}
For a modest explicit denominator, take $a=2799/1712$.
The following four pairs are the exact floors, with the minus sign first.
\begin{center}
\begin{tabular}{rrr}\toprule
$k$ & $\lfloor a^{3^k-k}\rfloor$ & $\lfloor a^{3^k+k}\rfloor$\\\midrule
1&2&7\\
2&31&223\\
3&133039&2540821\\
4&27506392487476241&1404192796017919199\\\bottomrule
\end{tabular}
\end{center}
All eight entries in the two output columns are prime. A second rational
base gives prime pairs at all seven levels $1\le k\le 7$.
The following PARI/GP code specifies this second base using integer roots
and thirteen small offsets. Each use of \texttt{sqrtnint} rounds down.
The fractions and floors are exact and \texttt{isprime} proves primality.
At $k=7$ the two primes have 387 and 389 decimal digits.
This finite computation makes no assertion about the next level.
{\small\begin{verbatim}
\\ Exact PARI/GP certificate for a common rational at seven levels.
expo=[2,4,7,11,24,30,77,85,238,248,723,735,2180,2194];
off=[1,1,4,15,4,104,16,75,12,192,108,2873,140];
p=vector(#expo); p[1]=2;
for(j=2,#expo,p[j]=sqrtnint(p[j-1]^expo[j],expo[j-1])+off[j-1]);
laste=expo[#expo]; den=10^401;
A=(1+sqrtnint(p[#p]*den^laste,laste))/den;
{for(j=1,#expo,
  if(floor(A^expo[j])!=p[j] || !isprime(p[j]), error("failure")))}
print("All fourteen floors are prime.");
\end{verbatim}}

\section{Prescribed decimal suffixes on fixed prime inputs}\label{sec:suffix}

Let $(P_k)$ be an increasing sequence of positive integers with
$\log P_{k+1}/\log P_k\to3$. Fix $0\le\delta<1/12$, and set
\[
 t_k=\lfloor\delta\log_{10}P_k\rfloor,\qquad M_k=10^{t_k}.
\]
Choose any residue $b_k$ coprime to $M_k$ at each index. Let $E$ be the
set of real $\lambda>1$ for which, at every sufficiently large $k$,
$\lfloor\lambda P_k\rfloor$ is a prime congruent to $b_k\bmod M_k$.

\begin{theorem}\label{thm:suffix-dimension}
For every nonempty open interval $I\subset(1,\infty)$,
\[
                 \dim_{\mathrm H}(E\cap I)=1-\frac{3\delta}{2}.
\]
The set $E$ has Lebesgue measure zero, including when $\delta=0$.
\end{theorem}

\begin{proof}
Work in a compact interval of positive length inside $I$. Write
$L_k=\log P_k$. The ratio assumption gives
$\sum_{i\le k}L_i=(3/2+o(1))L_k$, by geometric comparison, and
$\sum_{i\le k}\log L_i=o(L_k)$ after omitting finitely many initial terms.
At level $k$ retain compact cells
\[
 J_k(q)=\big[(q+1/4)/P_k,\,(q+3/4)/P_k\big]
\]
for admissible primes $q$. Given a parent at level $k-1$, the middle
half of its image under multiplication by $P_k$ has length
$H\asymp P_k/P_{k-1}=P_k^{2/3+o(1)}$, at height $x\asymp P_k$.
Every $q$ in this smaller interval gives a whole child inside its parent
for sufficiently large $k$. Uniformly in the parent and residue, the
number of such children is
\[
                   \gg \frac{P_k}{P_{k-1}M_k\log P_k}.
\]
Indeed, $H/\varphi(M_k)\ge x^{7/12+\eta}$ for a fixed $\eta>0$ since
$\delta<1/12$. For $\delta>0$ apply the short-interval prime number
theorem of Thorner and Zaman for moduli of fixed radical, in the form
\cite[Corollary 3.1 and (3.2) of the preprint]{tz}. For $\delta=0$ use
the ordinary short-interval prime number theorem in the same range.
The initial compact interval supplies a first cell by the same estimate
at a sufficiently large starting index.

Give each retained child an equal share of its parent's mass and write
$s=1-3\delta/2$. Telescoping the child bounds shows that every level-$k$
cell has mass at most
\[
  P_k^{-1}\exp\!\left(\sum_{i\le k}\log M_i+o(L_k)\right)
                         =P_k^{-s+o(1)} .
\]
For $P_k^{-1}\le h<P_{k-1}^{-1}$, an interval of length $h$ meets at most
$2+hP_k/M_k$ such cells, since all their labels belong to the same residue
class modulo $M_k$. Its mass is therefore at most
$2P_k^{-s+o(1)}+hP_k^{\delta/2+o(1)}$. For every $\varepsilon>0$ this
is $O_\varepsilon(h^{s-\varepsilon})$. The first term uses
$h\ge P_k^{-1}$ and the second uses
$\log(1/h)\ge L_{k-1}=(1/3+o(1))L_k$.
The mass distribution principle gives the lower bound.

For the upper bound use full floor cells of width $1/P_k$ and impose
only the congruences from any fixed starting index. A parent meets at
most $2+P_k/(P_{k-1}M_k)$ permitted cells at the next level. The product
of these bounds gives at most $P_k^{s+o(1)}$ cells in each fixed bounded
interval. Taking the countable union over starting indices and bounded
intervals proves the upper bound. Finally, on any fixed bounded interval
the union of all prime cells at level $k$ has length $O(1/\log P_k)$ by
the prime number theorem. The sum over $k$ converges. The first
Borel--Cantelli lemma proves the measure assertion.
\end{proof}

\begin{corollary}\label{cor:both-prime-histories}
Let $p_k^\pm=\lfloor A_\pm^{\,3^k\pm k}\rfloor$ be the primes of the two
prescribed minima and set $P_k=\max(p_k^+,p_k^-)$. In any nonempty open
interval $I\subset(1,\infty)$ there is a set of multipliers of Hausdorff
dimension $23/26$ such that, for every sufficiently large $k$,
$\lfloor\lambda P_k\rfloor$ is a prime whose decimal suffix contains,
in order, the complete decimal words for $p_{k-3}^+$ and $p_{k-3}^-$.
A set of the same dimension can instead record both entire prime
histories through index $k-4$ at every sufficiently large $k$. Those
histories can be read from the output prime alone.
\end{corollary}

\begin{proof}
Put $D_k=\log_{10}P_k$ and $L=\max(\log_{10} A_+,\log_{10} A_-)>0$.
Then $D_k=L3^k+O(k)$, so Theorem~\ref{thm:suffix-dimension} applies.
Use $\delta=1/13$. The combined digit count of $p_{k-3}^+$ and $p_{k-3}^-$
is at most $(2/27+o(1))D_k$, strictly below $D_k/13$. Concatenate
their words, end with the last prime digit, and pad the word on the left
with zeros to $t_k$ digits. Its last digit is coprime to ten for large
$k$. These are admissible prescribed residues, with finitely many
initial residues chosen arbitrarily.

For the second assertion use both words at each index $i\le k-4$.
Their combined digit count is at most
$2\sum_{i\le k-4}(L3^i+O(i))=(1/27+o(1))D_k$.
To make the histories recoverable from the prime alone, give each word
$w$ a length tag on its right. If $d$ is its length and $m$ is the
number of digits of $d$, form the tag consisting of $m$ copies of $1$,
a $0$, and the $m$ digits of $d$. Append this tag in reverse order to
$w$. Concatenate the tagged words chronologically, prefix a single $0$
as a stop symbol, and pad on the left. A reader moving from right to
left recovers $m$, then $d$, then the entire word, and repeats until
the stop symbol. The reader needs neither the input prime nor the
multiplier. The last digit is $1$, and the tags cost $O(k^2)$ digits
altogether. In both cases the
dimension with the specified padding is $1-3/(2\cdot13)=23/26$.
\end{proof}


\paragraph{Divisibility by earlier primes.}
A related consequence, with its own choice of $\lambda$, uses
multiplicative residues instead of decimal words. Put
$m(k)=\lfloor\log k/\log 3\rfloor$ and let $M'_k$ be the least common
multiple of all $p_i^+$ and $p_i^-$ with $1\le i\le m(k)$.
There is a set of $\lambda$ of Hausdorff dimension one in every open
$I\subset(1,\infty)$ such that, eventually,
\[
 Q_k=\lfloor\lambda P_k\rfloor\in\PP,\qquad M'_k\mid Q_k-1 .
\]
Thus each prime from either input sequence divides $Q_k-1$ at every
sufficiently late index. Indeed $\log M'_k=O(3^{m(k)})=O(k)$, which
by \eqref{eq:growth} gives $M'_k\le(\log P_k)^D$ for a fixed $D$.
The short-interval Siegel--Walfisz bound in \cite[p.~3 of the preprint]{tz}
is uniform for these moduli. The parent interval of length
$P_k^{2/3+o(1)}$ supplies $\gg P_k/(P_{k-1}M'_k\log P_k)$ children
with residue one. The proof of Theorem~\ref{thm:suffix-dimension} now
gives mass at most $P_k^{-1+o(1)}$ per cell, since
$\sum_{i\le k}\log M'_i=O(k^2)=o(\log P_k)$.
At length $P_k^{-1}\le h<P_{k-1}^{-1}$ at most $2+hP_k$ cells meet an
interval. This proves the asserted dimension by the same mass estimate.
\section{Prescribing the gap after each output prime}\label{sec:gaps}
For a prime $q$ let $g(q)$ be the distance from $q$ to the next prime.
Write $\log_j$ for the $j$-fold iterated natural logarithm and put
\[
  L(x)=\frac{\log x\,\log_2 x\,\log_4 x}{\log_3 x}
\]
for sufficiently large $x$.

\begin{theorem}
Let $(P_k)$ be a fixed sequence of primes with
$P_k\to\infty$ and $\log P_{k+1}/\log P_k\to3$.
There are absolute constants $D,c>0$ with the following property.
For any word $(w_k)\in\{\mathrm{s},\mathrm{l}\}^{\mathbb N}$
and any open interval $I\subset(0,\infty)$, there are an index
$K$ and a Cantor set of slopes $\lambda\in I$ such that, for every
$k\ge K$, $q_k=\lfloor\lambda P_k\rfloor$ is prime and
\[
 g(q_k)\le D\quad(w_k=\mathrm{s}),\qquad
 g(q_k)\ge cL(q_k)\quad(w_k=\mathrm{l}).
\]
In particular the theorem applies to the pointwise maximum
of two fixed prime sequences whose logarithms each have successive
ratio tending to $3$. A word using both letters infinitely often
gives $\liminf g(q_k)/\log q_k=0$ and
$\limsup g(q_k)/\log q_k=\infty$.
\end{theorem}

\begin{proof}
First consider any interval $[X,X+X^{3/5}]$ with $X$ sufficiently
large. Alweiss and Luo \cite{alweiss}, applied with exponent $0.55$
in a subinterval about its midpoint, give consecutive primes of
distance at most an absolute $D$. For the other gap size take
$u=\lfloor(\log X)/5\rfloor$ and $Q=\prod_{p\le u}p$.
Definition~1 and (1.1) of Ford, Green, Konyagin, Maynard and Tao
\cite{fgkmt} provide residues $a_p\bmod p$ for $p\le u$ that
cover every integer in $[1,H]$, where
$H=\lfloor c_0 u\log u\,\log_3 u/\log_2 u\rfloor$
for an absolute $c_0>0$. As $Q=X^{1/5+o(1)}$, the Chinese remainder
theorem supplies an integer $m$ in the middle third of our interval
with $m\equiv-a_p\pmod p$ for each $p\le u$. All integers
$m+1,\ldots,m+H$ are composite and $H\gg L(X)$.
There is a prime on each side of the run within the original
interval, by the same short-interval theorem of Alweiss and Luo
in the outer quarters. The consecutive primes bounding the run
thus lie in that interval, and the smaller one $q$ has
$g(q)\ge H+1\ge cL(q)$ for some absolute $c>0$.

Choose a compact interval of positive length inside $I$.
Its image under multiplication by $P_K$ contains an interval of the
above size with fixed margins when $K$ is large.
Choose a prime $q$ of the required type there and retain the cell
$[(q+1/4)/P_K,(q+3/4)/P_K]$. Given such a cell at level $k$,
choose two disjoint closed interior subintervals of length
$1/(10P_k)$ separated by at least $1/(10P_k)$.
Multiplication by $P=P_{k+1}$ sends either one to an interval
at scale $X\asymp_I P$ of length
$P/(10P_k)=P^{2/3+o(1)}$. Its middle half contains a full
$[X,X+X^{3/5}]$ and leaves more than one unit of endpoint clearance
in the original image. Each image therefore contains the whole
next cell $[q+1/4,q+3/4]$ at a prime of the prescribed type.
The estimates hold uniformly for all nodes once $K$ is sufficiently
large. The two disjoint children at every step and diameters
$1/(2P_k)\to0$ give a Cantor set on which all the stated
floor identities hold. For the assertion about maxima, apply the
inequalities $(3-\epsilon)\log P_k^{(j)}\le\log P_{k+1}^{(j)}
\le(3+\epsilon)\log P_k^{(j)}$ to both sequences and take
maxima. The two limit statements follow from
$L(x)/\log x\to\infty$.
\end{proof}

If $I$ contains $(1,1+\varepsilon)$ for some $\varepsilon>0$,
primality can be kept at every index while the prescribed gap word
starts at $K$. Indeed, take $N$ so large that $P_N$ dominates the
preceding inputs and start with the cell
$1+[1/(4P_N),3/(4P_N)]\subset I$. Every floor through $N$ is then
$P_j$ itself. The induction applies from $K=N+1$ onward.

\section{Polynomial values on three consecutive prime outputs}
For $s\in\{-1,1\}$ put $C_k=3^k+sk$. Given $A>1$, write
$p_k=\lfloor A^{C_k}\rfloor$. For a nonzero polynomial $F$ and a
positive tuple $\mathbf x$, let
$M_F(\mathbf x)=\max_{\nu\in\operatorname{supp} F}\mathbf x^\nu$.
The least prime-representing base for each of these clocks is
transcendental by \cite{corpus8}.

\begin{theorem}\label{thm:three-prime-separation}
The base $A$ is transcendental if and only if every nonzero
$F\in\overline{\mathbb Q}[X,Y,Z]$ satisfies, for some $c_F>0$ and all
sufficiently large $k$,
\[
 |F(p_k,p_{k+1},p_{k+2})|\ge c_F M_F(p_k,p_{k+1},p_{k+2}).
\]
In particular the inequality holds on the actual prime triples for
both minima. For four consecutive digits the corresponding assertion
fails for every $A>1$.
\end{theorem}
\begin{proof}
For a monomial with exponent vector $(i,j,\ell)$, put
$w=i+3j+9\ell$, $d=i+j+\ell$, and $h=j+2\ell$.
Since $p_m/A^{C_m}\to1$, its value is
$A^{w3^k+skd+sh}(1+o(1))$. Let $(w_*,v_*)$ be the lexicographically
largest pair $(w,sd)$ in the support. After division by
$A^{w_*3^k+kv_*}$, only monomials attaining this pair survive.
Their exponents $h$ are distinct because the three linear forms
$w,d,h$ are independent. The normalized sum therefore tends to the
value of a nonzero polynomial over $\overline{\mathbb Q}$ at $A^s$.
It is nonzero when $A$ is transcendental, while the normalized
$M_F$ has a positive finite limit. This proves the bound.

For the converse, the identity $3C_k+C_{k+2}-4C_{k+1}=-2s$ gives
\[
 \frac{p_k^3p_{k+2}}{p_{k+1}^4}\longrightarrow A^{-2s}.
\]
If $A$ is algebraic, let $P$ be the minimal polynomial of $A^{-2s}$
and $m=\deg P$. Then $F=Y^{4m}P(X^3Z/Y^4)$ is nonzero.
At the indicated triples, $F/Y^{4m}\to0$ and $M_F/Y^{4m}$ tends
to a positive number. Finally,
$3C_k+5C_{k+2}=7C_{k+1}+C_{k+3}$ gives the four-variable witness
$X^3Z^5-Y^7W$, whose two monomials have ratio tending to one.
\end{proof}

\section*{Scope}
In the common-base theorem the prime coordinate and the prime-free coordinate
are prescribed at every sufficiently large index. Whether one base can make
both coordinates prime for every $k\ge1$ remains a separate question.
The dilation constructions choose $\lambda$ after fixing the desired
conditions. Erd\H os's Problem~972 specifies an arbitrary irrational
$\lambda>1$ in advance \cite{erdos}. The results here therefore establish
existence statements with fixed input primes and additional output constraints.
They do not decide that universal problem. The paper gives no improvement
to the short-interval estimates it uses.

\begin{thebibliography}{9}
\bibitem{corpus8} Keio, \emph{Additive Clocks, Non-Reversible Recurrences, and
Transcendental Prime-Representing Constants}, manuscript 008, especially
the coefficient-family theorem and the final scope discussion, in the
\href{https://github.com/QEEDDD/Verifiable-Mathematical-Reasoning-Corpus/tree/ecd87de3348dcb706a76bb12fde722d7a348a012}{Verifiable Mathematical Reasoning Corpus, 2 November 2026 revision}.
\bibitem{corpus31} Keio, \emph{Prime-Representing Constants and Related Structures
(Encodings, Arithmetic, Geometry, and Modular Obstructions)}, manuscript
031 in the same corpus, section on logarithmic prime-free windows for
non-reversible clocks.
\bibitem{saito} K. Saito, \emph{Transcendency of variants of Mills' constant},
Ramanujan J. \textbf{70} (2026), article 69.
\href{https://arxiv.org/html/2508.16068v3}{arXiv:2508.16068v3}.
\bibitem{tz} J. Thorner and A. Zaman, \emph{Refinements to the prime number
theorem for arithmetic progressions}, Math. Z. \textbf{306} (2024).
\href{https://arxiv.org/pdf/2108.10878}{arXiv:2108.10878, Corollary 3.1 and (3.2)}.
\bibitem{watt} N. Watt, \emph{Short intervals almost all containing primes},
Acta Arith. \textbf{72} (1995), 131--167.
\href{https://doi.org/10.4064/aa-72-2-131-167}{doi:10.4064/aa-72-2-131-167}.
\bibitem{mv} H. L. Montgomery and R. C. Vaughan, \emph{The large sieve},
Mathematika \textbf{20} (1973), 119--134.
\href{https://doi.org/10.1112/S0025579300004708}{doi:10.1112/S0025579300004708}.
\bibitem{alweiss} R. Alweiss and S. Luo, \emph{Bounded gaps between primes in short
intervals}, Res. Number Theory \textbf{4} (2018).
\href{https://arxiv.org/abs/1707.05437}{arXiv:1707.05437}.
\bibitem{fgkmt} K. Ford, B. Green, S. Konyagin, J. Maynard and T. Tao,
\emph{Long gaps between primes}, J. Amer. Math. Soc. \textbf{31} (2018),
65--105. \href{https://doi.org/10.1090/jams/876}{doi:10.1090/jams/876}.
\bibitem{erdos} T. F. Bloom (ed.), \emph{Erd\H os Problems}, Problem 972,
attributed to Erd\H os (1965). \url{https://www.erdosproblems.com/972}.
\end{thebibliography}
\end{document}
